% ===== wrap-up =====
\begin{frame}{이번 주차 정리}
  \small
  \begin{enumerate}
    \item \kb{3.1 베이즈 정리 = ``같은 $P(x,y)$ 를 앞뒤로 나누는''}
      \hbox{}\\
      세 단어(사전·가능도·사후) + ``세로 합 = 조건부 세계'' 확률
      트리. \textbf{오즈 = 가능도비 곱의 연쇄}가 NB의 로그 합과
      로지스틱 시그모이드를 같은 자리에서 설명. \textbf{기저율
      오류}: 사전확률을 빼먹으면 99\% 검진도 9\% 사후확률.
    \item \kb{3.2 GDA = 가우시안 + 공유 공분산 + MLE}
      \hbox{}\\
      평균·공분산으로 \textbf{닫힌 형태}. 이차항 상쇄 $\Rightarrow$
      \textbf{결정 경계는 항상 직선}. $P(y{=}1|x)=\sigma(w^\top x+b)$
      --- \textbf{로지스틱회귀와 같은 시그모이드}. $\phi$ 는
      위치만 미는 \textbf{병렬 이동} = 경계의 기저율 오류.
    \item \kb{3.3 NB = 독립 가정 + 빈도 카운팅}
      \hbox{}\\
      파라미터 \textbf{선형}($2V{+}1$) vs GDA \textbf{제곱}($V^2{+}V$)
      $\Rightarrow$ \textbf{차원의 저주 회피}. 틀려도 \textbf{동일한
      방향}으로 틀려 순위 유지. 스무딩(0 강제 방지) · \textbf{로그
      합}(언더플로우) · \textbf{사전확률}(가능도비 2도 10\%면 뒤집힘).
  \end{enumerate}
  \vskip0.8em
  \begin{block}{다음 주 --- Week 4. kNN과 k-means}
    또 다른 극단: \textbf{배울 파라미터가 전부 없는} 분류기.
    ``$k$개 가장 가까운 이웃을 보고 다수결'' --- ``모델 없이
    거리만으로''. 이번 장의 ``데이터 생성 법칙을 배운다''에서
    ``아무것도 배우지 않는다''로 넘어가는 것. 다만 이번 장이 예고한
    \textbf{차원의 저주}는 Week 4에서 ``왜 kNN이 고차원에서 무너지는가''
    로 다시 나옴.
  \end{block}
\end{frame}
